\documentclass[10pt]{article}
\usepackage[T1]{fontenc}\usepackage{lmodern}\pdfmapfile{+lm.map}
\usepackage[margin=0.8in]{geometry}
\usepackage{amsmath,amssymb,booktabs,longtable,array,graphicx,xcolor,listings,hyperref}
\definecolor{navy}{HTML}{193959}\definecolor{green}{HTML}{19745C}
\hypersetup{colorlinks=true,linkcolor=navy,urlcolor=navy}
\setlength{\parindent}{0pt}\setlength{\parskip}{6pt}
\newcommand{\ps}[2]{\langle #1\rangle\,\psi #2}
\newcommand{\same}[1]{\left.#1\right|_{\rm hand}=\left.#1\right|_{\rm compiler}}
\lstset{language=C,basicstyle=\ttfamily\scriptsize,lineskip=-0.6pt,breaklines=true,columns=fullflexible,keepspaces=true,showstringspaces=false,numbers=left,numberstyle=\tiny,frame=single}

\title{Hand-to-Compiler Correspondence for Papers I--IV\\Final DNFs, ONFs, Lifted ONFs and C}
\author{Lenore M. Mullin\\College of Nanotechnology, Science and Engineering\\University at Albany\\\texttt{lmullin@albany.edu}}
\date{Completed correspondence study; consolidated release, 24 November 2026}
\begin{document}\maketitle
\textbf{Conclusion.} The recovered hand C implementation inventory is closed.
Every numerical routine/region in the 20 pinned hand C files has a compiler
counterpart. All 39 archived regions pass, including the original untiled/tiled
routines, every numerical serial reference, and all fixed and benchmark variants.
Paper III adds its printed OpenACC kernel and four explicitly derived C ONFs of
its grouped-head/cache Python algorithms, giving 44 paired cases.

All 246 independently recovered end-DNF component comparisons pass. All 44 pairs
have matching addressed operational C ASTs and loop/directive structures. Their
lifted ONFs are equal for a common admissible machine shape and mapping. All
138 untimed finite comparisons give zero hand/compiler difference. There are no
remaining equivalence rows in this declared C-emission study.

\begin{center}\begin{tabular}{lr}\toprule
Check & Result\\\midrule
Archived numerical C regions / source files & 39 / 20, all covered\\
Total paired cases & 44 / 44 pass\\
Recovered DNF components & 246 / 246 equal\\
Matched operational and lifted structures & 44 / 44 equal\\
Untimed finite comparisons & 138 / 138 pass\\
Matched mathematical helper comparisons & 27 / 27 pass\\
Deliberately corrupted C ASTs rejected & 109 / 109\\
Additional mathematical defects rejected & 15 / 15\\
Owner / tile / grouped-head regressions & 259 / 280 / 119 pass\\
Corrected paired fresh-device entry contracts & 5 / 5 pass\\\bottomrule
\end{tabular}\end{center}

The equality is \emph{within each hand/compiler pair}. Different algorithms and
different hand schedules need not share one DNF or one ONF. Derivation steps may
differ; the final checked forms agree. This is a restricted, trusted executable
translation-validation result, with the premises stated below. It does not
claim a proof-assistant-verified universal compiler or new hardware timing.

\textbf{Documents.} This report states the result and its proof boundary. The
complete code volume prints all 20 original hand source files, all 88 paired
kernel translation units, and ten corrected entry-contract units. The AST atlas
shows all 44 control/ownership summaries. Full expression trees, all 246 formula
pairs, code, figures and replay tools are in the accompanying package.

\newpage\section{Scope and provenance: nothing left unpaired}
The numerical inventory is derived from function definitions in every archived
C file. Each function is classified as a certified numerical region, a matched
mathematical helper, or input/output, timing and verification support. The tool
fails on an unclassified numerical function. All 20 files and all 39 numerical
regions are covered; the seven previously unpaired source files are resolved.

The archived C files are byte-identical to published repository commit
\texttt{0cf1d1f} (full identity in the evidence). Original files are preserved
unchanged. Common test adapters expose the numerical region as a kernel API;
I/O and timing drivers are not asserted byte-identical to those adapters.

Paper III requires an explicit source distinction. Its C/OpenACC Listing 3,
printed page 8, was manually transcribed and checked against the supplied PDF.
The grouped GQA/MQA algorithm is published as Python Listing 7, and the
preallocated cache as Python Listing 5. Their C reference ONFs were independently
written from those equations/listings and compared with compiler emission.
They are \emph{new hand C reference derivations}, not fabricated archived C files.
The source-PDF digest and transcription provenance are supplied. Grouped and
cache references retain the published default FP32 contract; historical decode
FP32 and FP64 interfaces remain distinct.

The recovered Fortran comparison programs and PyTorch autograd scripts provide
historical/oracle context. The closed claim here is C emission for the hand
algorithm families and contracts; it is not a verified Fortran or PyTorch frontend.

\subsection{Contracts}
The contracts require legal shapes, contiguous stated layouts, successful
allocations, valid nonaliasing where required, representable integer address
intermediates, standard attention scale $1/\sqrt D$, and RMSNorm
$\epsilon=10^{-6}$. Training C interfaces use the common feature extent $D$.
Grouped attention requires $G\mid H$. Cache length satisfies $0\le N<C$,
including the initial empty cache. Tiling requires $T>0$ and $T\mid N$.
Fixed scratch bounds are enforced as premises: $N\le8192$ where \texttt{MAX\_N}
is used, and the relevant $D,D_{\rm FF}\le512$ where 512-element arrays are indexed.
Additive targets require zero entry states where the routine does not assign
those identities. Five separate-device profiles have matched corrected data
contracts that supply those identities.

DNF equality is in the exact-real abstraction of the typed operations. The C
comparison retains floating types and scalar operation order. It establishes
the same admissible floating-point behavior under common runtime choices, not
uniform error bounds against real arithmetic. Numerical checks are finite:
paired error is zero; the largest independent NumPy-oracle difference is about
$1.40\times10^{-8}$ in the FP32 cases.

\newpage\section{The independently checked denotations}
The following equalities summarize the final hand and compiler DNFs. Scalar
abbreviations below stand for substituted expressions, not allocated intermediate
arrays. All fully expanded/canonical component pairs and digests are supplied.
For attention, define
\[
 s_{bik}=D^{-1/2}\sum_{j=0}^{D-1}(\ps{b,i,j}{Q})(\ps{b,k,j}{K}),
 \qquad a_{bik}=\frac{e^{s_{bik}}}{\sum_{r=0}^{N-1}e^{s_{bir}}}.
\]
The forward output is the same final expression on both paths:
\[
 \ps{b,i,d}{Out_{\rm hand}}=\ps{b,i,d}{Out_{\rm compiler}}
 =\sum_{k=0}^{N-1}a_{bik}(\ps{b,k,d}{V}).
\]
For stored-$A$ backward set $w_{bik}=\ps{b,i,k}{A}$; for fused
forward/backward set $w_{bik}=a_{bik}$. Define
\[
 g_{bik}=\sum_{d=0}^{D-1}(\ps{b,i,d}{dO})(\ps{b,k,d}{V}),\quad
 c_{bi}=\sum_{r=0}^{N-1}w_{bir}g_{bir},\quad
 z_{bik}=D^{-1/2}w_{bik}(g_{bik}-c_{bi}).
\]
Both paths recover
\begin{align*}
\ps{b,k,d}{GV}&=\sum_i w_{bik}(\ps{b,i,d}{dO}),\\
\ps{b,i,j}{GQ}&=\sum_k z_{bik}(\ps{b,k,j}{K}),\\
\ps{b,k,j}{GK}&=\sum_i z_{bik}(\ps{b,i,j}{Q}).
\end{align*}
Every sum has the legal zero-based bound shown in its component record. The
untiled originals, tiled originals, atomic and atomic-free variants recover
these same algorithm-specific denotations, while preserving their distinct ONFs.

For grouped decode, $g(h)=\lfloor h/(H/G)\rfloor$, and
\[
 u_{hk}=D^{-1/2}\sum_j(\ps{h,j}{Q})(\ps{g(h),k,j}{K}),\qquad
 \ps{h,d}{Out}=\sum_k\frac{e^{u_{hk}}}{\sum_r e^{u_{hr}}}(\ps{g(h),k,d}{V}).
\]
The same expression covers GQA, MQA ($G=1$), MHA ($G=H$), and the singleton
query-head specialization. The separate rank-one interface keeps $\rho q=\langle D\rangle$;
it introduces no artificial matrix category for a vector.

Cache append writes only
\[
\ps{g,N,j}{K'}=\ps{g,j}{key},\quad
\ps{g,N,d}{V'}=\ps{g,d}{value},\qquad N'=N+1.
\]
All other cache components are unchanged. The proof checks every store, not just
the final stored rule, so a hidden prefix write is rejected. The empty-cache case
is included. Capacity remains preallocated; no existing prefix is copied.

\newpage\section{Normalization, FFN and complete-block closure}
Use token extent $L=BN$ and $\rho X=\langle L,D\rangle$. Define
\[
 r_t=\sqrt{D^{-1}\sum_j(\ps{t,j}{X})^2+\epsilon},\quad
 \widehat x_{td}=(\ps{t,d}{X})/r_t,\quad
 c_t=\sum_j(\ps{t,j}{GZ})(\ps{j}{\Gamma})\widehat x_{tj}.
\]
The hand and compiler normalization components agree:
\begin{align*}
\ps{t,d}{Z}&=(\ps{d}{\Gamma})\widehat x_{td},\\
\ps{t,d}{GX}&=\bigl[(\ps{t,d}{GZ})(\ps{d}{\Gamma})-\widehat x_{td}c_t/D\bigr]/r_t,\\
\ps{d}{G\Gamma}&=\sum_t(\ps{t,d}{GZ})\widehat x_{td}.
\end{align*}
For gated FFN, with hidden extent $F$, let
\[
 u_{tf}=\sum_j(\ps{t,j}{X})(\ps{j,f}{Wg}),\quad
 v_{tf}=\sum_j(\ps{t,j}{X})(\ps{j,f}{Wu}),\quad h_{tf}=\operatorname{silu}(u_{tf})v_{tf}.
\]
Then $\ps{t,d}{F_{\rm out}}=\sum_f h_{tf}(\ps{f,d}{Wd})$ on both paths.
Writing $\eta_{tf}=\sum_d(\ps{t,d}{GF})(\ps{f,d}{Wd})$, the recovered
local derivatives are $g^v_{tf}=\eta_{tf}\operatorname{silu}(u_{tf})$ and
$g^u_{tf}=\eta_{tf}v_{tf}\operatorname{silu}'(u_{tf})$. Thus both paths give
\begin{align*}
\ps{t,d}{GX}&=\sum_f\bigl[g^u_{tf}(\ps{d,f}{Wg})+g^v_{tf}(\ps{d,f}{Wu})\bigr],\\
\ps{d,f}{GWg}&=\sum_t(\ps{t,d}{X})g^u_{tf},\qquad
\ps{d,f}{GWu}=\sum_t(\ps{t,d}{X})g^v_{tf},\\
\ps{f,d}{GWd}&=\sum_t h_{tf}(\ps{t,d}{GF}).
\end{align*}
SiLU and its derivative are expanded algebraically in the component evidence;
the matched C helper bodies are checked separately.

The complete-block proof is compositional through named producer/consumer cuts.
Each component is checked before substitution into its consumers. There are
three 30-node complete-block graphs and a 31-node fixed serial graph; each closes
all ten roots: $Y$, $GX$, two scale gradients and six weight gradients. Two
10-node norm--FFN-chain graphs close eight roots each. Extra serial stages and
retained intermediates are preserved rather than silently fused away.

The original fixed serial complete block has a separate attention output and
residual, separate FFN projection/output stages, and separate residual-gradient
seeds. The compiler now reproduces those choices too. Matching the final DNF
never required those derivation steps to coincide with another implementation's.

\newpage\section{ONFs, gamma addresses and dimension lifting}
For each pair the comparison retains typed declarations, scalar arithmetic order,
loops and bounds, local allocation/free positions, data clauses, helpers,
reductions, atomics and parallel directives. Normalization removes formatting and
single-statement loop braces, canonicalizes exact decimal spelling, and compares
integer address polynomials under the no-overflow premise. It does not
reassociate floating-point sums/products. This proves equality of computational
C ASTs; entire I/O/timing drivers need not be byte-identical.

For the original tile lift $N=N_TT$,
\[
 i=it\,T+tt,\quad 0\le it<N_T,\quad0\le tt<T,
 \qquad (it,tt)=(\lfloor i/T\rfloor,i\bmod T).
\]
Consequently $\rho K_{\rm lifted}=\langle B,N_T,T,D\rangle$ and
$\tau(K_{\rm lifted})=BN_TTD=BND=\tau(K)$. The addresses agree:
\[
\gamma(\langle b,it,tt,d\rangle,\langle B,N_T,T,D\rangle)
 =(bN+itT+tt)D+d.
\]
The interpreter recovers affine tile writes through this inverse and collapses
the resulting nested finite reductions through the bijection. Tile-stride,
tile-extent and accumulation defects are rejected. This is a symbolic
reindexing argument supported by 280 finite tile checks.

For machine owner count $P=\tau(M)$ and phase work cardinality $\ell>0$,
let $a_r=\lfloor\ell r/P\rfloor$, $b_r=\lfloor\ell(r+1)/P\rfloor$.
The owner/local-index relation $i=a_r+j$, $0\le j<b_r-a_r$, has inverse
$r=\lceil P(i+1)/\ell\rceil-1$, $j=i-a_r$. Its count telescopes to $\ell$.
If $P\mid\ell$, the rectangular lift has shape
$\rho M\mathbin{+\!+}\langle\ell/P\rangle$ and the same component count.
Otherwise the certificate uses a segmented legal-index relation; padded storage
capacity is not incorrectly equated to the original logical $\tau$.

Every loop and directive is independently extracted from both C regions,
including Listing 3's nested data/vector constructs. The ordered structures
agree in all 44 cases. Hence, for the same admissible machine shape, binding
$\mathcal B$, workspace assignment and runtime choices,
\[
\Omega_{M,\mathcal B}(ONF_{\rm hand})=
\Omega_{M,\mathcal B}(ONF_{\rm compiler}).
\]
This is equality of a parameterized family. Serial references retain their serial
binding; a parallel extension requires an admissible private-workspace/reduction
mapping. No claim of identical independently chosen runtime placements follows.
The ten-field machine-property order remains
$\langle U_d,W_d,M_d,t_d,nmax_d,Mcap_d,P_d,X_d,C_d,R_d\rangle$.
Physical values remain parameters rather than invented measurements.

\newpage\section{Entry repairs and what can be concluded without experiments}
Five archival GPU profiles need a corrected fresh-device initialization contract:
the benchmark and fixed complete block, both norm--FFN chain variants, and the
historical fused-attention atomic variant. Host-zeroed additive targets were
listed under ordinary \texttt{create}/\texttt{copyout}; those clauses do not
initialize fresh separate device storage from host zeros. A focused static audit
identifies the first additive writes and confirms the affected targets.

Both hand and compiler paths now have separately saved corrected variants using
\texttt{copy} for those host-zeroed accumulators. The arithmetic bodies are
unchanged. All five corrected pairs have equal operational kernel ASTs and pass
syntax checks. Original archival source/data contracts remain untouched. The
DNF proof's zero-entry premise is therefore explicit and realizable under the
corrected fresh-allocation contract.

The corrections change transfer accounting. For a full block with $L=BN$, FP64
initial copies add $8(2LD+2D+3D^2+3DF)$ bytes; moving $GQ,GK$ from
\texttt{create} to \texttt{copy} additionally adds their final $16LD$-byte copy.
For a norm--FFN chain the added initial payload is $8(D+3DF)$ bytes. The atomic
fused-attention variant adds the initial $8BND$ bytes for $GK$. These are logical
data-clause payloads under fresh separate allocations, not measured DRAM traffic.
No new device execution is claimed.

\subsection{Connection to Paper V and Paper I}
Paper V's hand-realization measurements and diagnostics remain the empirical
foundation. The present result establishes compiler correspondence to the pinned
hand implementations and documented contracts. Its published source snapshot
matches all 20 recovered C files. The Paper I A100 job 137615 archive also passes
all 403 frozen-file/runner/compiler binding checks. Those observations keep their
original implementation and machine identities.

Some older Paper V jobs lack an immutable executed-source digest and build
identity in the recovered records. Atomic and atomic-free revisions differ. Those historical
measurements are not relabeled as executions of the current reference compiler
or the corrected entry profiles. Quantitative transfer requires the measured
source, precision, outputs, machine, mapping and measurement boundary to match.
The source correspondence supports carrying the mathematical operation and
logical-access analysis across a matched realization. The cost analysis must
retain each version's actual AROW/GS storage and atomic updates. A DNF alone does not
supply an absolute runtime, energy, bandwidth or occupancy prediction.

\textbf{Study conclusion.} The hand-to-compiler normal-form/code correspondence
study is complete for the recovered C implementation inventory and the stated
Paper III source contracts. New timing experiments are unnecessary to establish
these equalities. Any later device/performance experiments test machine/runtime
assumptions or new empirical claims; they are not an outstanding equivalence gate.
This result closes the reproduction study while leaving the broader predictive
performance research program free to proceed on its own evidence.

\newpage\section{Final coverage ledger}
Every row below has matching DNF components, addressed ONF/C structure and
common-machine lifted structure. ``Arch.'' denotes a pinned archived C region,
``Print'' the transcribed printed C listing, and ``Deriv.'' an explicitly derived
C ONF of a published Python algorithm. ``Repair'' means that a matched corrected
fresh-device entry pair is included alongside the archival pair.
\small\begin{longtable}{p{2.9in}lrlp{.65in}}\toprule
Case & Origin & DNFs & Entry & Status\\\midrule\endhead
\texttt{p4\_\allowbreak norm\_\allowbreak ffn\_\allowbreak gpu} & Arch. & 10 & Repair & Pass\\
\texttt{p3\_\allowbreak decode\_\allowbreak cpu} & Arch. & 1 & Std. & Pass\\
\texttt{p4\_\allowbreak block\_\allowbreak cpu} & Arch. & 30 & Std. & Pass\\
\texttt{p4\_\allowbreak block\_\allowbreak gpu} & Arch. & 30 & Repair & Pass\\
\texttt{p4\_\allowbreak rmsnorm\_\allowbreak cpu} & Arch. & 3 & Std. & Pass\\
\texttt{p4\_\allowbreak ffn\_\allowbreak cpu} & Arch. & 5 & Std. & Pass\\
\texttt{p1\_\allowbreak forward\_\allowbreak cpu} & Arch. & 2 & Std. & Pass\\
\texttt{p1\_\allowbreak forward\_\allowbreak gpu} & Arch. & 2 & Std. & Pass\\
\texttt{p2\_\allowbreak backward\_\allowbreak cpu} & Arch. & 3 & Std. & Pass\\
\texttt{p2\_\allowbreak backward\_\allowbreak gpu} & Arch. & 3 & Std. & Pass\\
\texttt{p2\_\allowbreak fused\_\allowbreak cpu} & Arch. & 4 & Std. & Pass\\
\texttt{p2\_\allowbreak fused\_\allowbreak gpu} & Arch. & 4 & Std. & Pass\\
\texttt{p3\_\allowbreak decode\_\allowbreak gpu} & Arch. & 1 & Std. & Pass\\
\texttt{p4\_\allowbreak rmsnorm\_\allowbreak verify} & Arch. & 3 & Std. & Pass\\
\texttt{p4\_\allowbreak ffn\_\allowbreak verify} & Arch. & 5 & Std. & Pass\\
\texttt{p4\_\allowbreak block\_\allowbreak verify\_\allowbreak cpu} & Arch. & 31 & Std. & Pass\\
\texttt{p4\_\allowbreak block\_\allowbreak verify\_\allowbreak gpu} & Arch. & 30 & Repair & Pass\\
\texttt{p4\_\allowbreak chain\_\allowbreak bench\_\allowbreak gpu} & Arch. & 10 & Repair & Pass\\
\texttt{p2\_\allowbreak fused\_\allowbreak atomic\_\allowbreak gpu} & Arch. & 4 & Repair & Pass\\
\texttt{original\_\allowbreak forward\_\allowbreak dnf} & Arch. & 2 & Std. & Pass\\
\texttt{original\_\allowbreak forward\_\allowbreak onf} & Arch. & 2 & Std. & Pass\\
\texttt{original\_\allowbreak backward\_\allowbreak dnf} & Arch. & 3 & Std. & Pass\\
\texttt{original\_\allowbreak backward\_\allowbreak onf} & Arch. & 3 & Std. & Pass\\
\texttt{original\_\allowbreak fused\_\allowbreak dnf} & Arch. & 4 & Std. & Pass\\
\texttt{original\_\allowbreak fused\_\allowbreak onf} & Arch. & 4 & Std. & Pass\\
\texttt{ref\_\allowbreak forward\_\allowbreak cpu} & Arch. & 2 & Std. & Pass\\
\texttt{ref\_\allowbreak forward\_\allowbreak gpu} & Arch. & 2 & Std. & Pass\\
\texttt{ref\_\allowbreak backward\_\allowbreak forward} & Arch. & 2 & Std. & Pass\\
\texttt{ref\_\allowbreak backward\_\allowbreak serial} & Arch. & 3 & Std. & Pass\\
\texttt{ref\_\allowbreak backward\_\allowbreak gpu\_\allowbreak forward} & Arch. & 1 & Std. & Pass\\
\texttt{ref\_\allowbreak backward\_\allowbreak gpu\_\allowbreak serial} & Arch. & 3 & Std. & Pass\\
\texttt{ref\_\allowbreak fused\_\allowbreak forward} & Arch. & 2 & Std. & Pass\\
\texttt{ref\_\allowbreak fused\_\allowbreak backward} & Arch. & 3 & Std. & Pass\\
\texttt{ref\_\allowbreak fused\_\allowbreak gpu\_\allowbreak forward} & Arch. & 2 & Std. & Pass\\
\texttt{ref\_\allowbreak fused\_\allowbreak gpu\_\allowbreak backward} & Arch. & 3 & Std. & Pass\\
\texttt{ref\_\allowbreak fused\_\allowbreak v3\_\allowbreak forward} & Arch. & 2 & Std. & Pass\\
\texttt{ref\_\allowbreak fused\_\allowbreak v3\_\allowbreak backward} & Arch. & 3 & Std. & Pass\\
\texttt{ref\_\allowbreak decode\_\allowbreak double} & Arch. & 1 & Std. & Pass\\
\texttt{ref\_\allowbreak decode\_\allowbreak float} & Arch. & 1 & Std. & Pass\\
\texttt{paper3\_\allowbreak listing3} & Print & 1 & Std. & Pass\\
\texttt{paper3\_\allowbreak grouped\_\allowbreak openmp} & Deriv. & 3 & Std. & Pass\\
\texttt{paper3\_\allowbreak grouped\_\allowbreak openacc} & Deriv. & 3 & Std. & Pass\\
\texttt{paper3\_\allowbreak append\_\allowbreak openmp} & Deriv. & 5 & Std. & Pass\\
\texttt{paper3\_\allowbreak append\_\allowbreak openacc} & Deriv. & 5 & Std. & Pass\\
\bottomrule\end{longtable}\normalsize
\section{Reproducibility and final files}
Run \texttt{replay\_all.sh} to regenerate and check all pairs, formula/graph
records, helper matches, deliberate defects, lifting certificates, coverage
classification and the final seal. It builds untimed local kernels, explicitly
using OpenACC host fallback. It submits no job and records no performance.
Run \texttt{build.sh} to compile this report, the complete code volume and AST
atlas. The latter's diagrams summarize control and ownership; full arithmetic
and expression ASTs are supplied as text/JSON for every case.

The trusted basis is visible: independently entered hand equations and source
transcriptions, the restricted recipe interpreter and reindexing rules, SymPy
exact algebra, helper matching, C parsing and the stated execution contracts.
The result is an executable, fail-closed correspondence certificate for this
implementation family, not a universal proof about all future compiler inputs.

\nocite{paper1,paper2,paper3,paper4,paper5,openacc}

\section{Consolidation: preserve the completed result}
The 44-pair study above is already complete. The November 24 release brings its existing emitter, independent checkers, pinned hand sources and certificates together with the separate compiler-audit branch. A shared command exports each attention case's existing DNF, addressed ONF, common-machine lifted record and generated C after checking recipe and operational-AST identity. These are integration checks, not a new derivation study. Corrected entry profiles are the default where fresh device storage needs explicit zero-state transfer.

The newer bounded Fortran extension has its own structured arithmetic and versioned-state schema. Its serial, OpenMP, OpenACC host and MPI numerical checks do not replace the attention evidence. In particular, lexicographic and snapshot-island sweeps have different finite-sweep state semantics. A common command interface is not yet a universal parser or a fully shared cross-family AST.

The persistent release ledger marks completed work separately from remaining language, target and predictive-model extensions. Paper I's historical physical-device evidence and Paper V's original measurements retain their own implementation/build/machine identities. No new GPU timing is added by consolidating these sources.

\bibliographystyle{plain}
\bibliography{references}
\par Published hand source: \href{https://github.com/womenflyplanes/moa-attention-verified-mullin}{pinned implementation repository}.
Compiler and replay checks: \href{https://github.com/womenflyplanes/moa-compiler/tree/codex/papers-i-iv-correspondence}{completed correspondence branch}.

\textbf{Acknowledgment.} University at Albany computational resources and support
are acknowledged. This correspondence replay ran no new work on UA or ACCESS.
\end{document}
